\documentclass[11pt]{article}

\usepackage[T1]{fontenc}
\usepackage[utf8]{inputenc}
\usepackage{lmodern}
\usepackage[margin=1in]{geometry}
\usepackage{microtype}
\usepackage{amsmath,amssymb,amsthm,mathtools}
\usepackage{stmaryrd}
\usepackage{enumitem}
\usepackage[hidelinks]{hyperref}

\newtheorem{theorem}{Theorem}[section]
\newtheorem{lemma}[theorem]{Lemma}
\newtheorem{proposition}[theorem]{Proposition}
\newtheorem{remark}[theorem]{Remark}

\title{A Note on Reciprocal-Chart Folding over Quartic Extensions}
\author{Seongjun Park\\Independent Researcher\\\texttt{parksj0106@gmail.com}}
\date{}

\begin{document}

\maketitle

\begin{abstract}
We study reciprocal-chart folding over a quartic extension \(E/F\). Fix a seed \(\tau\in E\) with \(F(\tau)=E\), and let the public orbit be defined by
\[
r_1=\tau,
\qquad
r_{i+1}=\frac{1}{r_i-c_i},
\qquad
c_i=
\begin{cases}
0,& i\text{ odd},\\
1,& i\text{ even}.
\end{cases}
\]
For \(F\)-valued leaves, we show that every round admits an exact reciprocal transport in the moving basis \(\{1,r_i,r_i^2,r_i^3\}\). This yields, for every depth \(n\ge 1\), a seed-specialized public \(F\)-linear evaluator
\[
T_{\tau,n}:F^{2^n}\to F^4
\]
whose intermediate values are always represented by vectors in \(F^4\).

Our main theorem gives an exact descriptor for this evaluator family. From the seed descriptor \((n,\tau)\), one can compute in \(O(n)\) base-field operations a reduced fold descriptor of size exactly \(4(n-1)\) field words together with \(8\) field words of reconstruction data. The reduced descriptor supports \(O(n)\)-time coefficient queries for entries of the matrix of \(T_{\tau,n}\), and an explicit evaluation algorithm computes \(T_{\tau,n}(x)\) in \(O(N)\) time for \(N=2^n\). Under the canonical straight-line implementation analyzed here, the algorithm performs exactly \(2N-4\) base-field multiplications on \(F\)-valued leaves.

This paper formalizes a cryptographically relevant algebraic backend for downstream commitment/opening protocols and non-interactive proof compilers: a constant-width seed-specialized evaluator, its exact descriptor size, and the structural implication for downstream interfaces that a layer restricted to explicit row vectors cannot preserve the \(O(n)\)-word specialization. A minimal information-theoretic oracle-model one-shot soundness theorem for the evaluator relation, with the output fixed before scalarization and with fresh verifier randomness, is also established.
\end{abstract}

\section{Introduction}

Reciprocal-chart folding appears naturally in sumcheck-style folding questions \cite{LFKN92,BootleChiesaSotiraki2021}. Once a post-commit seed \(\tau\) is fixed, however, the most natural abstraction is not an interactive transcript but the induced linear evaluator
\[
T_{\tau,n}:F^{2^n}\to F^4,
\]
obtained by iterating a public M\"obius orbit and an exact round-by-round reciprocal transport. The purpose of this paper is to formalize that evaluator family in the first nontrivial closed-form case: a quartic extension \(E/F\) with the alternating \(0/1\) orbit.

Existing presentations of sumcheck-style folding \cite{LFKN92,BootleChiesaSotiraki2021} and related algebraic backends typically describe the round-by-round update at a schematic level, leaving the concrete transport identities, descriptor sizes, and arithmetic costs implicit. In particular, for extensions of degree greater than two, no prior work provides an exact closed-form local fold identity in a moving power basis, a fully explicit reduced descriptor with a provable size bound, or the precise multiplication count of the resulting algorithm. This gap matters for downstream layers: without an exact algebraic specification, commitment/opening protocols and non-interactive compilers must either re-derive the round structure internally or fall back to generic interfaces that forfeit the logarithmic-size seed specialization.

This backend viewpoint is cryptographically useful for two complementary reasons. First, commitment and opening layers for structured linear families must decide whether they can consume a seed-dependent structured descriptor or whether they must materialize explicit rows \cite{CampanelliEtAl2022LMVC,CatalanoFioreTucker2022AHFC,BalbasCatalanoFioreLai2022CFC,FischLiuVesely2024Orbweaver}. Second, later non-interactive or correlated-challenge analyses are cleaner once the algebraic backend has already been stated as a precise linear object \cite{CanettiEtAl2018FS,BlockEtAl2023FRI}. In the quartic reciprocal chart studied here, that backend is highly structured: each active intermediate value remains in \(F^4\), the seed specialization is describable in \(O(n)\) words, and explicit-row interfaces lose that succinctness.

\paragraph{Contributions.}
Our main contributions are as follows.
\begin{enumerate}[label=\arabic*.]
\item We prove an \emph{exact local reciprocal fold identity} in a moving quartic basis.
\item We prove a \emph{structured evaluator theorem}: for \(F\)-valued leaves, the seed-specialized map \(T_{\tau,n}\) admits a reduced public descriptor of size exactly \(4(n-1)\), an \(O(n)\)-time coefficient oracle, and an \(O(N)\)-time evaluation procedure.
\item We prove an \emph{exact canonical multiplication count} of \(2N-4\) for the explicit evaluation algorithm applied to \(F\)-valued leaves.
\item We establish an \emph{implication for downstream interfaces}: any downstream layer that insists on explicit row vectors must materialize the four rows of \(T_{\tau,n}\), and therefore cannot preserve the \(O(n)\)-word seed-dependent specialization.
\end{enumerate}

\subsection*{Related work and positioning}

Classical sumcheck originates with Lund, Fortnow, Karloff, and Nisan \cite{LFKN92}; modern cryptographic presentations include \cite{BootleChiesaSotiraki2021}, and recent implementation-oriented analyses include \cite{BagadDombThaler2024SmallChar,Rothblum2024MLENote}. On the commitment side, succinct openings for structured linear families are studied by, among others, \cite{CampanelliEtAl2022LMVC,CatalanoFioreTucker2022AHFC,BalbasCatalanoFioreLai2022CFC,FischLiuVesely2024Orbweaver}. On the folding side, recent protocol-level systems such as Nova, ProtoStar, and HyperNova develop recursive or accumulation-based proof layers above algebraic backends \cite{KothapalliSettyTzialla2021Nova,BuenzChen2023ProtoStar,KothapalliSetty2023HyperNova}. This paper presents a seed-specialized linear evaluator and establishes a minimal one-shot soundness statement in the oracle model that later cryptographic layers may refine.

\subsection*{Scope and organization}

We focus on the quartic case \([E:F]=4\) and on the alternating \(0/1\) schedule because this is the first degree in which the moving basis, reciprocal transport, coefficient update, and exact arithmetic constants all fit in a short self-contained treatment. The aim is precision rather than maximal generality. Sections~2--4 develop the algebraic backend, Section~5 gives the worked \(N=4\) example, Section~6 fixes the arithmetic cost model, Section~7 establishes the implication for explicit-row interfaces, and Section~8 offers concluding remarks. The appendix establishes a minimal information-theoretic soundness statement in an idealized oracle model where the claimed output is fixed before scalarization; concrete commitment/opening and Fiat--Shamir layers build on this foundation as subsequent work.

\section{Algebraic Setup and Reciprocal Orbit}

Sections~2 through~5 develop the algebraic backend, with Section~5 presenting the first detailed example. Section~6 then defines an arithmetic cost model in order to state exact multiplication counts. Let \(F\) be a field of characteristic different from \(2\) and \(3\), let \(E/F\) be a quartic extension, and let
\[
N=2^n,
\qquad n\ge 1.
\]
We index leaves by
\[
j\in\{0,\dots,N-1\},
\qquad
j=\sum_{i=1}^{n} b_i2^{i-1},
\qquad
b_i\in\{0,1\},
\]
using the least-significant-first path convention.

Whenever an integer \(k\) appears in a field expression, it denotes its image in \(F\).

Define the algebraic seed domain
\[
\Omega^{\mathrm{alg}}:=\{\tau\in E : F(\tau)=E\}.
\]

For \(\tau\in\Omega^{\mathrm{alg}}\), define the public challenge orbit by
\[
r_1=\tau,
\qquad
r_{i+1}=\frac{1}{r_i-c_i},
\qquad
c_i=
\begin{cases}
0,& i\text{ odd},\\
1,& i\text{ even}.
\end{cases}
\]

\begin{lemma}[Orbit formulas and moving charts]
\label{lem:orbit}
Let \(\tau\in\Omega^{\mathrm{alg}}\). Then every reciprocal update is well-defined. Moreover, for every admissible \(k\) one has
\[
r_{2k+1}=\frac{\tau}{1-k\tau},
\qquad
r_{2k+2}=\frac{1-k\tau}{\tau},
\]
and for every admissible round \(i\),
\[
F(r_i)=E.
\]
In particular, \(\{1,r_i,r_i^2,r_i^3\}\) is an \(F\)-basis of \(E\) at every round.
\end{lemma}

\begin{proof}
Since \(\tau\notin F\), one has \(\tau\neq 0\). If \(1-k\tau=0\) for some admissible integer \(k\), then the image of \(k\) in \(F\) is nonzero, for otherwise the left-hand side would equal \(1\). Hence \(k\) is invertible in \(F\), so \(\tau=k^{-1}\in F\), contradicting \(F(\tau)=E\). Therefore, the denominators in the displayed closed forms are nonzero.

The formulas follow by induction on \(k\). The base cases \(r_1=\tau\) and \(r_2=\tau^{-1}\) are immediate. If \(r_{2k+1}=\tau/(1-k\tau)\), then
\[
r_{2k+2}=\frac{1}{r_{2k+1}}=\frac{1-k\tau}{\tau}.
\]
If \(r_{2k+2}=(1-k\tau)/\tau\), then
\[
r_{2k+3}=\frac{1}{r_{2k+2}-1}=\frac{1}{(1-(k+1)\tau)/\tau}=\frac{\tau}{1-(k+1)\tau}.
\]
Each \(r_i\) is obtained from \(\tau\) by an invertible M\"obius transformation with coefficients in \(F\), so \(F(r_i)=F(\tau)=E\). Since \([E:F]=4\), the displayed set is an \(F\)-basis.
\end{proof}

For each round \(i\), let
\[
m_i(X)=X^4-\mu_{i,3}X^3-\mu_{i,2}X^2-\mu_{i,1}X-\mu_{i,0}
\]
be the monic minimal polynomial of \(r_i\) over \(F\), so that
\[
r_i^4=\mu_{i,3}r_i^3+\mu_{i,2}r_i^2+\mu_{i,1}r_i+\mu_{i,0}.
\]
Write
\[
\mu_i:=(\mu_{i,0},\mu_{i,1},\mu_{i,2},\mu_{i,3})\in F^4.
\]

An active entry at round \(i\) is stored as a four-word state
\[
s=(a_0,a_1,a_2,a_3)\in F^4,
\]
which, together with a shared denominator \(\Delta_i\in E^\times\), represents the element
\[
\llbracket s\rrbracket_i
=
\Delta_i^{-1}\big(a_0+a_1r_i+a_2r_i^2+a_3r_i^3\big).
\]
Equivalently, if
\[
P_s(Y):=a_0+a_1Y+a_2Y^2+a_3Y^3,
\]
then \(\llbracket s\rrbracket_i=\Delta_i^{-1}P_s(r_i)\).

For a scalar input \(v\), the initial state is
\[
\Delta_1=1,
\qquad
s(v)=(v,0,0,0)\quad(v\in F).
\]
The shared denominator evolves by
\[
\Delta_{i+1}=\Delta_i r_{i+1}^3.
\]

\begin{lemma}[Closed forms for the shared denominator]
\label{lem:delta}
For every admissible integer \(k\), one has
\[
\Delta_{2k}=\tau^{-3},
\qquad
\Delta_{2k+1}=(1-k\tau)^{-3}.
\]
\end{lemma}

\begin{proof}
The base cases are \(\Delta_1=1=(1-0\cdot\tau)^{-3}\) and
\[
\Delta_2=\Delta_1 r_2^3=\tau^{-3}.
\]
Assume \(\Delta_{2k}=\tau^{-3}\). Using Lemma~\ref{lem:orbit},
\[
\Delta_{2k+1}=\Delta_{2k}r_{2k+1}^3
=\tau^{-3}\left(\frac{\tau}{1-k\tau}\right)^3
=(1-k\tau)^{-3}.
\]
Then
\[
\Delta_{2k+2}=\Delta_{2k+1}r_{2k+2}^3
=(1-k\tau)^{-3}\left(\frac{1-k\tau}{\tau}\right)^3
=\tau^{-3}.
\]
This proves both formulas by induction.
\end{proof}

\section{Exact Local Reciprocal Fold}

Fix an admissible round \(i\), and write \(r=r_i\), \(c=c_i\), and \(\mu=\mu_i\). For adjacent states
\[
s_0=(a_0,a_1,a_2,a_3),
\qquad
s_1=(b_0,b_1,b_2,b_3),
\]
set \(\delta_j=b_j-a_j\). Define the intermediate state \(u=(u_0,u_1,u_2,u_3)\) by
\[
u_0=a_0+\mu_0\delta_3,
\qquad
u_1=a_1+\delta_0+\mu_1\delta_3,
\]
\[
u_2=a_2+\delta_1+\mu_2\delta_3,
\qquad
u_3=a_3+\delta_2+\mu_3\delta_3.
\]
Equivalently,
\[
u=s_0+S_{\mu}(s_1-s_0),
\qquad
S_{\mu}=
\begin{pmatrix}
0&0&0&\mu_0\\
1&0&0&\mu_1\\
0&1&0&\mu_2\\
0&0&1&\mu_3
\end{pmatrix}.
\]
The chart transport map is
\[
R_0(u_0,u_1,u_2,u_3)=(u_3,u_2,u_1,u_0),
\]
for odd rounds, and
\[
R_1(u_0,u_1,u_2,u_3)=
\big(u_3,\ u_2+3u_3,\ u_1+2u_2+3u_3,\ u_0+u_1+u_2+u_3\big),
\]
for even rounds.

The Step-A update satisfies the identity
\[
P_u(r)=P_{s_0}(r)+r\big(P_{s_1}(r)-P_{s_0}(r)\big).
\]
In particular, all seed dependence in the local fold is completely captured by the coefficient tuple \(\mu_i\), and within Step A it affects the computation only through the single scalar \(\delta_3\).

\begin{lemma}[Transport identities]
\label{lem:transport}
For every \(u\in F^4\),
\[
P_{R_0u}(Y)=Y^3P_u(Y^{-1}),
\qquad
P_{R_1u}(Y)=Y^3P_u(1+Y^{-1}).
\]
Consequently, if \(r_{i+1}=1/(r_i-c_i)\) and \(\Delta_{i+1}=\Delta_i r_{i+1}^3\), then for \(s'=R_{c_i}u\),
\[
\Delta_{i+1}^{-1}P_{s'}(r_{i+1})
=
\Delta_i^{-1}P_u(r_i).
\]
\end{lemma}

\begin{proof}
The displayed polynomial identities are immediate from the definitions of \(R_0\) and \(R_1\). For \(c_i=0\), one has \(r_i=r_{i+1}^{-1}\); for \(c_i=1\), one has \(r_i=1+r_{i+1}^{-1}\). Substituting these relations and using \(\Delta_{i+1}=\Delta_i r_{i+1}^3\) yields the claim.
\end{proof}

\begin{theorem}[Exact reciprocal-chart fold]
\label{thm:exact-fold}
Let \(u=s_0+S_{\mu_i}(s_1-s_0)\) and \(s'=R_{c_i}u\). Then
\[
\llbracket s'\rrbracket_{i+1}
=
\llbracket s_0\rrbracket_i+r_i\big(\llbracket s_1\rrbracket_i-\llbracket s_0\rrbracket_i\big).
\]
In particular, each fold step from length \(N/2^{i-1}\) to length \(N/2^i\) is exact, the active state entries remain in \(F^4\), and the denominator update is globally shared.
\end{theorem}

\begin{proof}
By the quartic relation for \(r_i\),
\[
P_u(r_i)=P_{s_0}(r_i)+r_i\big(P_{s_1}(r_i)-P_{s_0}(r_i)\big).
\]
Applying Lemma~\ref{lem:transport} to \(s'=R_{c_i}u\) gives
\[
\Delta_{i+1}^{-1}P_{s'}(r_{i+1})=\Delta_i^{-1}P_u(r_i),
\]
which is exactly the displayed identity after substituting the definition of \(P_u(r_i)\).
\end{proof}

\section{Global Structured Evaluator and Reduced Descriptor}

The local rule of Theorem~\ref{thm:exact-fold} is linear over \(F\). At round \(i\), each adjacent pair is updated by
\[
(s_{2t},s_{2t+1})\longmapsto A_i s_{2t}+B_i s_{2t+1},
\]
where
\[
A_i=R_{c_i}(I-S_{\mu_i}),
\qquad
B_i=R_{c_i}S_{\mu_i}.
\]

The first round simplifies further because the leaves are scalar states.

\begin{lemma}[Seed-independence of the first round]
\label{lem:first-round}
For any \(v_0,v_1\in F\),
\[
R_0\big(s(v_0)+S_{\mu_1}(s(v_1)-s(v_0))\big)=(0,0,v_1-v_0,v_0),
\]
and this value is independent of \(\mu_1\).
\end{lemma}

\begin{proof}
Since \(s(v_1)-s(v_0)=(v_1-v_0,0,0,0)\), one has
\[
S_{\mu_1}(s(v_1)-s(v_0))=(0,v_1-v_0,0,0).
\]
Thus
\[
s(v_0)+S_{\mu_1}(s(v_1)-s(v_0))=(v_0,v_1-v_0,0,0),
\]
and applying \(R_0\) yields the claim.
\end{proof}

For all complexity analyses, we work in the finite-field specialization
\[
F=\mathbb F_q,
\qquad
E=\mathbb F_{q^4},
\qquad
\operatorname{char}(F)>3,
\]
and count arithmetic operations over \(F\) in the unit-cost model. Public descriptors are encoded relative to a fixed ambient \(F\)-basis of \(E\), while intermediate fold states continue to use the moving round basis \(\{1,r_i,r_i^2,r_i^3\}\).

Define the public descriptors
\[
d^{\mathrm{seed}}_{\tau,n}:=(n,\tau),
\qquad
d^{\mathrm{fold}}_{\tau,n}:=(\mu_2,\dots,\mu_n)\in F^{4(n-1)},
\]
where for \(n=1\) the fold descriptor is the empty tuple, and
\[
d^{\mathrm{rec}}_{\tau,n}:=(r_{n+1},\Delta_{n+1})\in E\times E^\times.
\]
The fold descriptor describes the evaluator on \(F\)-valued leaves; the reconstruction descriptor is only needed to recover the terminal \(E\)-valued quantity represented by the final state.

For the seed-independent first round, define the fixed branch vectors
\[
\nu^{(0)}:=(0,0,-1,1)^\top,
\qquad
\nu^{(1)}:=(0,0,1,0)^\top.
\]
For rounds \(i\ge 2\), define the branch matrices
\[
M_i^{(0)}:=A_i,
\qquad
M_i^{(1)}:=B_i.
\]

\begin{theorem}[Structured evaluator and reduced descriptor]
\label{thm:descriptor}
For every \(\tau\in\Omega^{\mathrm{alg}}\), the reciprocal folding construction defines a public \(F\)-linear map
\[
T_{\tau,n}:F^N\to F^4.
\]
Moreover:
\begin{enumerate}[label=\arabic*.]
\item There is a deterministic expansion algorithm
   \[
   \mathsf{Expand}:(n,\tau)\longmapsto \big(d^{\mathrm{fold}}_{\tau,n},d^{\mathrm{rec}}_{\tau,n}\big)
   \]
   running in \(O(n)\) base-field operations.
\item The reduced fold descriptor has size exactly \(4(n-1)\) base-field words. Relative to a fixed ambient \(F\)-basis of \(E\), the reconstruction descriptor has size exactly \(8\) base-field words.
\item For each \(j=\sum_{i=1}^{n} b_i2^{i-1}\), the \(j\)-th column of \(T_{\tau,n}\) is
   \[
   T_{\tau,n}(e_j^{(N)})=
   \begin{cases}
   \nu^{(b_1)},& n=1,\\[0.5ex]
   M_n^{(b_n)}\cdots M_2^{(b_2)}\nu^{(b_1)},& n\ge 2,
   \end{cases}
   \]
   where \(e_j^{(N)}\) is the \(j\)-th standard basis vector of \(F^N\). Hence any coefficient of the matrix of \(T_{\tau,n}\) is computable in \(O(n)\) time from \(d^{\mathrm{fold}}_{\tau,n}\).
\item Given \(d^{\mathrm{fold}}_{\tau,n}\) and \(x\in F^N\), the explicit evaluation algorithm computes \(T_{\tau,n}(x)\) in \(O(N)\) base-field operations using \(O(1)\) field words per active array entry.
\end{enumerate}
\end{theorem}

\begin{proof}
Linearity follows because each round is an explicit \(F\)-linear block map and the initialization is linear. For expansion, represent \(E\) relative to a fixed ambient \(F\)-basis and compute the minimal polynomial of \(\tau\) using constant-size linear algebra; this costs \(O(1)\) because the extension degree is fixed at \(4\). Proposition~\ref{prop:coeff-update} gives an \(O(1)\)-time update for the coefficient tuple at each round, and the orbit point \(r_{i+1}\) and shared denominator \(\Delta_{i+1}\) are also updated in \(O(1)\) time per round. Thus \(\mathsf{Expand}\) runs in \(O(n)\) time. The output includes only \(\mu_2,\dots,\mu_n\) because Lemma~\ref{lem:first-round} shows that the first round on scalar leaves is seed-independent. Item (2) is immediate from the descriptor definitions.

For Item (3), the first path bit \(b_1\) selects one of the two fixed round-1 branch vectors \(\nu^{(0)}\) or \(\nu^{(1)}\). Each subsequent bit \(b_i\) selects the left or right branch matrix \(M_i^{(0)}\) or \(M_i^{(1)}\). Composing these choices gives the displayed column formula. Each step is a constant-size matrix-vector multiplication, so the cost is \(O(n)\). Item (4) is exactly the explicit evaluation algorithm of Theorem~\ref{thm:exact-fold} applied layer by layer to the state array.
\end{proof}

\begin{remark}[Reduced versus uniform per-round descriptors]
\label{rem:tight-vs-uniform}
The theorem above states the reduced descriptor size for the evaluator \(T_{\tau,n}:F^N\to F^4\) on \(F\)-valued leaves. If one instead wants a uniform round-by-round transcript, one may also store \(\mu_1\), obtaining a descriptor of size exactly \(4n\). This extra first-round tuple does not affect the algebraic structure and is not needed for the evaluator itself.
\end{remark}

In particular, there exist unique row vectors
\[
\lambda_{\tau,n}^{(0)},\lambda_{\tau,n}^{(1)},\lambda_{\tau,n}^{(2)},\lambda_{\tau,n}^{(3)}\in F^N
\]
such that
\[
T_{\tau,n}(x)=
\big(
\langle\lambda_{\tau,n}^{(0)},x\rangle,
\langle\lambda_{\tau,n}^{(1)},x\rangle,
\langle\lambda_{\tau,n}^{(2)},x\rangle,
\langle\lambda_{\tau,n}^{(3)},x\rangle
\big).
\]
The terminal extension-field value represented by the final state is recovered by the public reconstruction map
\[
\mathsf{Rec}_{\tau,n}(z_0,z_1,z_2,z_3)
:=
\Delta_{n+1}^{-1}\big(z_0+z_1r_{n+1}+z_2r_{n+1}^2+z_3r_{n+1}^3\big)\in E.
\]

\begin{remark}[Scope of the quartic case]
\label{rem:quartic}
The quartic case is singled out because it already yields a nontrivial constant-width evaluator with fully explicit local transport, a closed-form coefficient update, and exact arithmetic constants. Higher-degree analogues are a natural next direction.
\end{remark}

\begin{remark}[Materializing explicit rows]
\label{rem:rows}
The full \(4\times N\) matrix of \(T_{\tau,n}\) can be explicitly formed in \(O(N)\) time by reverse propagation through the depth-\(n\) binary circuit. This is routine once the branch vectors and branch matrices are known, but it matters later when determining whether a subsequent protocol can take the structured descriptor directly as input, or whether it requires explicit rows.
\end{remark}

\section[Worked Example: N=4]{Worked Example: \(N=4\)}

The case \(N=4=2^2\) is the first nontrivial instance and is worth writing out explicitly. Let
\[
x=(x_0,x_1,x_2,x_3)\in F^4.
\]
At round \(1\), the two leaf pairs are \((x_0,x_1)\) and \((x_2,x_3)\). By Lemma~\ref{lem:first-round}, the first round is seed-independent:
\[
s_{01}^{(1)}=(0,0,x_1-x_0,x_0),
\qquad
s_{23}^{(1)}=(0,0,x_3-x_2,x_2).
\]
In particular, \(\mu_1\) does not appear anywhere in the evaluator for \(N=4\).

Now let \(\mu_2=(\mu_{2,0},\mu_{2,1},\mu_{2,2},\mu_{2,3})\). The single surviving pair at round \(2\) is
\[
s_0=(0,0,x_1-x_0,x_0),
\qquad
s_1=(0,0,x_3-x_2,x_2).
\]
Hence
\[
\delta_0=\delta_1=0,
\qquad
\delta_2=(x_3-x_2)-(x_1-x_0),
\qquad
\delta_3=x_2-x_0,
\]
and Step A yields
\[
u_0=\mu_{2,0}(x_2-x_0),
\qquad
u_1=\mu_{2,1}(x_2-x_0),
\]
\[
u_2=(x_1-x_0)+\mu_{2,2}(x_2-x_0),
\qquad
u_3=x_0+\delta_2+\mu_{2,3}(x_2-x_0).
\]
Applying the even-round transport gives the final state
\[
T_{\tau,2}(x)=R_1u=(z_0,z_1,z_2,z_3),
\]
with explicit matrix form
\[
T_{\tau,2}(x)=
\begin{pmatrix}
2-\mu_{2,3} & -1 & \mu_{2,3}-1 & 1\\
5-\mu_{2,2}-3\mu_{2,3} & -2 & \mu_{2,2}+3\mu_{2,3}-3 & 3\\
4-\mu_{2,1}-2\mu_{2,2}-3\mu_{2,3} & -1 & \mu_{2,1}+2\mu_{2,2}+3\mu_{2,3}-3 & 3\\
1-\sum_{j=0}^{3}\mu_{2,j} & 0 & \sum_{j=0}^{3}\mu_{2,j}-1 & 1
\end{pmatrix}
\begin{pmatrix}
x_0\\x_1\\x_2\\x_3
\end{pmatrix}.
\]
This makes the global linearity explicit: the four displayed rows are precisely
\[
\lambda_{\tau,2}^{(0)},\lambda_{\tau,2}^{(1)},\lambda_{\tau,2}^{(2)},\lambda_{\tau,2}^{(3)}.
\]

The multiplication count also matches the general theorem of the next section: round \(1\) requires zero multiplications, and round \(2\) requires exactly four, hence
\[
M_{\mathrm{canon}}(4)=4=2\cdot 4-4.
\]

\section{Exact Canonical Multiplication Count}

We determine the exact multiplication complexity of the explicit evaluation algorithm under a simple straight-line implementation model.
\begin{itemize}[leftmargin=1.5em]
\item Leaves are initialized as scalars \((v,0,0,0)\).
\item In round \(1\), scalar inputs are processed by the specialized routine
  \[
  \mathsf{FoldPairScalar}(v_0,v_1):=(0,0,v_1-v_0,v_0),
  \]
  which is exactly the odd-round transport of \((v_0,v_1-v_0,0,0)\) and uses no field multiplication.
\item From round \(2\) onward, the algorithm uses the generic Step-A update with the four products
  \[
  \mu_0\delta_3,\quad \mu_1\delta_3,\quad \mu_2\delta_3,\quad \mu_3\delta_3,
  \]
  followed by the public transport \(R_0\) or \(R_1\).
\item Additions, subtractions, rewiring, and add-chain realizations of multiplication by the public constants \(2\) and \(3\) incur no multiplicative cost.
\item No further runtime sparsity specialization is performed after round \(1\).
\end{itemize}

\begin{proposition}[Exact canonical multiplication count]
\label{prop:canon}
Assume \(N=2^n\) with \(n\ge 1\). Under the canonical implementation above, the explicit evaluation algorithm performs exactly
\[
M_{\mathrm{canon}}(N)=2N-4
\]
base-field multiplications on \(F\)-valued leaves.

Here the count refers only to the local fold arithmetic once the public tuples \(\mu_2,\dots,\mu_n\) are already known. It excludes descriptor expansion, terminal reconstruction in \(E\), and any commitment/opening overhead.
\end{proposition}

\begin{proof}
Round \(1\) contains \(N/2\) pairs, and each is processed by \(\mathsf{FoldPairScalar}\), so round \(1\) requires zero multiplications. Every later pair requires exactly four multiplications from Step~A and none from transport. The number of pairs after round \(1\) is
\[
\sum_{i=2}^{n}\frac{N}{2^i}=\frac{N}{2}-1.
\]
Therefore
\[
M_{\mathrm{canon}}(N)=4\left(\frac{N}{2}-1\right)=2N-4.
\]
\end{proof}

If each multiplication by the public constants \(2\) and \(3\) is counted as a full field multiplication, then every even-round transport \(R_1\) can be realized with at most two additional multiplications per surviving pair: one for \(3u_3\), and one for \(2u_2\), while \(2u_3=(3u_3)-u_3\) is then computed without additional multiplications. Summing those additional costs over even rounds gives
\[
M_{\mathrm{cons}}(N)
\le
M_{\mathrm{canon}}(N)+2\sum_{\substack{2\le i\le n\\ i\text{ even}}}\frac{N}{2^i}
< 2N-4+\frac{2N}{3}
=\frac{8}{3}N-4.
\]

Proposition~\ref{prop:canon} captures a structural property of the reciprocal chart: all seed dependence is concentrated into the single scalar \(\delta_3\) per surviving pair, while the chart transport itself is multiplication-free under the canonical model.

\subsection*{Direct Fixed-Basis Baseline}

To establish a baseline for comparison, consider the most direct ambient-basis implementation of the same logical fold: each intermediate value is stored in a fixed \(F\)-basis of \(E\), and each surviving pair is updated by
\[
w' = w_0 + r_i(w_1-w_0),
\]
where \(r_i\in E\) is the public round challenge. Let \(m_{F,\mathrm{pub}}\) denote the number of base-field multiplications needed to compute \(r\delta\) when \(r\in E\) is a generic public constant and \(\delta\in F\), and let \(m_{E,\mathrm{pub}}\) denote the corresponding count when \(\delta\in E\), under the chosen quartic multiplication routine.

\begin{proposition}[Direct fixed-basis baseline]
\label{prop:baseline}
Under the same multiplication-dominant model as Proposition~\ref{prop:canon}, the direct fixed-basis update performs
\[
M_{\mathrm{dir}}(N)
=
\frac{N}{2}\,m_{F,\mathrm{pub}}
+
\left(\frac{N}{2}-1\right)m_{E,\mathrm{pub}}
\]
base-field multiplications on \(F\)-valued leaves. Consequently,
\[
\frac{M_{\mathrm{dir}}(N)}{M_{\mathrm{canon}}(N)}
\longrightarrow
\frac{m_{F,\mathrm{pub}}+m_{E,\mathrm{pub}}}{4}
\qquad\text{as }N\to\infty.
\]
\end{proposition}

\begin{proof}
In round \(1\), each pair begins from \(F\)-valued leaves, so \(w_1-w_0\in F\) and the direct update requires \(m_{F,\mathrm{pub}}\) multiplications per pair. There are \(N/2\) such pairs. From round \(2\) onward, each current state is a general element of \(E\), so each surviving pair incurs \(m_{E,\mathrm{pub}}\) multiplications, and the number of such pairs is
\[
\sum_{i=2}^{n}\frac{N}{2^i}=\frac{N}{2}-1.
\]
Summing the two contributions gives the formula, and the asymptotic ratio follows from Proposition~\ref{prop:canon}.
\end{proof}

\begin{remark}[Concrete baseline constants]
The symbolic parameters \(m_{F,\mathrm{pub}}\) and \(m_{E,\mathrm{pub}}\) encapsulate the choice of ambient basis and quartic multiplication routine. Proposition~\ref{prop:baseline} provides the general formulation; any concrete numerical instantiation should be stated relative to an explicitly fixed ambient representation.
\end{remark}

\section{Implication for Explicit-Row Interfaces}

Theorems~\ref{thm:exact-fold} and \ref{thm:descriptor} identify a concrete backend object for later cryptographic layers: the seed-specialized evaluator \(T_{\tau,n}\) together with its reduced descriptor \(d^{\mathrm{fold}}_{\tau,n}\). A subsequent commitment/opening protocol has two standard interface choices. It can take the structured descriptor directly as input, as in systems designed for structured linear families \cite{CampanelliEtAl2022LMVC,CatalanoFioreTucker2022AHFC,BalbasCatalanoFioreLai2022CFC,FischLiuVesely2024Orbweaver}, or it can explicitly form the four rows of \(T_{\tau,n}\) and use those rows.

The next proposition formalizes the cost of the row-materialization approach.

\begin{proposition}[Explicit-row interfaces preclude descriptor succinctness]
\label{prop:explicit-row}
Fix \(\tau\in\Omega^{\mathrm{alg}}\). Let a subsequent opening protocol express the \(k\)-th output coordinate of \(T_{\tau,n}\) as an explicit row vector
\[
w_{\tau,n}^{(k)}\in F^N
\]
satisfying
\[
\langle w_{\tau,n}^{(k)},x\rangle=\big(T_{\tau,n}(x)\big)_k
\qquad
\text{for all }x\in F^N.
\]
Then necessarily
\[
w_{\tau,n}^{(k)}=\lambda_{\tau,n}^{(k)}.
\]
Hence any explicit-row specialization to \(\tau\) requires exactly \(4N\) base-field coefficients, whereas the reduced fold descriptor has size \(4(n-1)\) (or \(4n\) in the uniform per-round encoding of Remark~\ref{rem:tight-vs-uniform}).
\end{proposition}

\begin{proof}
By definition of the row vectors \(\lambda_{\tau,n}^{(k)}\), we have
\[
\langle \lambda_{\tau,n}^{(k)},x\rangle=\big(T_{\tau,n}(x)\big)_k
\qquad
\forall x\in F^N.
\]
If \(w_{\tau,n}^{(k)}\) is another row vector satisfying the same identity, then
\[
\langle w_{\tau,n}^{(k)}-\lambda_{\tau,n}^{(k)},x\rangle=0
\qquad
\forall x\in F^N.
\]
Taking \(x=e_j^{(N)}\) for each standard basis vector shows
\[
w_{\tau,n}^{(k)}[j]=\lambda_{\tau,n}^{(k)}[j]
\qquad
\forall j,
\]
so the two rows are equal.
\end{proof}

Proposition~\ref{prop:explicit-row} is an interface statement, not a hardness statement. By Remark~\ref{rem:rows}, the rows can be explicitly formed in linear time from the structured circuit of Theorem~\ref{thm:descriptor}. A row-only interface inevitably replaces a logarithmic-size seed specialization by \(4N\) explicit coefficients. Therefore, any subsequent protocol that wants to preserve the \(O(n)\)-word specialization must take the structured reciprocal descriptor itself, or an equivalent structured representation, as input rather than only explicit row vectors.

\section{Concluding Remarks}

From the perspective of protocol design, the separation established in Proposition~\ref{prop:explicit-row} provides a clear interface contract. A commitment/opening mechanism can exploit the evaluator family's exact width, exact descriptor size, and exact fold arithmetic \cite{CampanelliEtAl2022LMVC,CatalanoFioreTucker2022AHFC,BalbasCatalanoFioreLai2022CFC,FischLiuVesely2024Orbweaver}. Likewise, a Fiat--Shamir or correlated-challenge analysis can treat the reciprocal chart as a fixed algebraic primitive rather than re-deriving its local identities within the security proof \cite{CanettiEtAl2018FS,BlockEtAl2023FRI}. The appendix provides a minimal information-theoretic soundness statement for the evaluator relation in the oracle model.

\appendix

\section{Reciprocal-Shift Coefficient Update}

\begin{proposition}[Reciprocal-shift minimal-polynomial update]
\label{prop:coeff-update}
Let
\[
m_i(X)=X^4-\mu_{i,3}X^3-\mu_{i,2}X^2-\mu_{i,1}X-\mu_{i,0}
\]
be the monic minimal polynomial of \(r_i\), fix \(c\in\{0,1\}\), and define
\[
D_c(\mu_i):=m_i(c)=c^4-\mu_{i,3}c^3-\mu_{i,2}c^2-\mu_{i,1}c-\mu_{i,0}.
\]
Then \(D_c(\mu_i)\neq 0\), and the next orbit point
\[
r_{i+1}=\frac{1}{r_i-c}
\]
has monic minimal polynomial
\[
m_{i+1}(Y)=\frac{Y^4m_i(c+Y^{-1})}{D_c(\mu_i)}.
\]
Equivalently,
\[
m_{i+1}(Y)=Y^4-\mu_{i+1,3}Y^3-\mu_{i+1,2}Y^2-\mu_{i+1,1}Y-\mu_{i+1,0},
\]
where
\[
\mu_{i+1,0}=-D_c(\mu_i)^{-1},
\qquad
\mu_{i+1,1}=(\mu_{i,3}-4c)D_c(\mu_i)^{-1},
\]
\[
\mu_{i+1,2}=(\mu_{i,2}+3c\mu_{i,3}-6c^2)D_c(\mu_i)^{-1},
\]
\[
\mu_{i+1,3}=(\mu_{i,1}+2c\mu_{i,2}+3c^2\mu_{i,3}-4c^3)D_c(\mu_i)^{-1}.
\]
\end{proposition}

\begin{proof}
Because \(m_i\) is irreducible of degree \(4\) over \(F\), it has no root in \(F\). Hence \(D_c(\mu_i)=m_i(c)\neq 0\) for every \(c\in F\), in particular for \(c\in\{0,1\}\).

Define
\[
\widetilde m(Y):=\frac{Y^4m_i(c+Y^{-1})}{D_c(\mu_i)}\in F[Y].
\]
Substituting \(Y=r_{i+1}\) gives
\[
\widetilde m(r_{i+1})
=
\frac{r_{i+1}^4 m_i(c+r_{i+1}^{-1})}{D_c(\mu_i)}
=
\frac{r_{i+1}^4 m_i(r_i)}{D_c(\mu_i)}
=0,
\]
so \(\widetilde m\) annihilates \(r_{i+1}\). Moreover, the coefficient of \(Y^4\) in \(Y^4m_i(c+Y^{-1})\) is precisely \(m_i(c)=D_c(\mu_i)\), so \(\widetilde m\) is a monic polynomial of degree \(4\).

The M\"obius relations
\[
r_i=c+r_{i+1}^{-1},
\qquad
r_{i+1}=(r_i-c)^{-1}
\]
show that \(r_i\in F(r_{i+1})\) and \(r_{i+1}\in F(r_i)\). Therefore
\[
F(r_{i+1})=F(r_i)=E,
\]
so the minimal polynomial of \(r_{i+1}\) over \(F\) has degree \(4\). Since \(\widetilde m\) is already a monic degree-\(4\) polynomial annihilating \(r_{i+1}\), it must equal the minimal polynomial \(m_{i+1}\).

It remains to expand. A direct computation gives
\[
Y^4m_i(c+Y^{-1})
=
D_c(\mu_i)Y^4
+
\bigl(4c^3-3c^2\mu_{i,3}-2c\mu_{i,2}-\mu_{i,1}\bigr)Y^3
\]
\[
\qquad
+
\bigl(6c^2-3c\mu_{i,3}-\mu_{i,2}\bigr)Y^2
+
\bigl(4c-\mu_{i,3}\bigr)Y
+1.
\]
Dividing by \(D_c(\mu_i)\) and comparing with the monic quartic form
\[
Y^4-\mu_{i+1,3}Y^3-\mu_{i+1,2}Y^2-\mu_{i+1,1}Y-\mu_{i+1,0}
\]
yields exactly the displayed coefficient formulas.
\end{proof}

\section{Information-Theoretic Oracle-Model Soundness}
\label{sec:it-soundness}

This section establishes an information-theoretic soundness theorem for the evaluator relation
\[
\mathcal R_{\tau,n}:=\{(x,y)\in F^N\times F^4 : y=T_{\tau,n}(x)\}.
\]
The main result is the following:
\begin{quote}
Suppose that (i) the prover's claimed output $y\in F^4$ is fixed before a random scalarization vector $\rho\in F^4$ is drawn, (ii) the sumcheck phase uses fresh, independent verifier challenges in $F$, and (iii) the verifier has exact access to $\widetilde x(z)$ at the final sampled point $z$. Then a false relation claim $y\neq T_{\tau,n}(x)$ is accepted with probability at most $(2n+1)/q$.
\end{quote}
The same section also explains why the more naive deterministic-orbit transcript is completely unsound. The error bound is nontrivial in the natural regime $q\gg n$; its leading $1/q$ term reflects the unavoidable collision probability of a single random scalarization.

\subsection{Setup}

Throughout this section, the protocol-level probability bounds rely only on the finite-field assumptions
\[
F=\mathbb F_q,
\qquad
N=2^n,
\quad n\ge 1.
\]
To instantiate the specific reciprocal evaluator of this paper, we further assume
\[
E=\mathbb F_{q^4},
\qquad
\operatorname{char}(F)>3,
\]
fix a seed $\tau\in E$ with $F(\tau)=E$, and let
\[
r_1=\tau,
\qquad
r_{i+1}=\frac{1}{r_i-c_i},
\qquad
c_i=
\begin{cases}
0,& i\text{ odd},\\
1,& i\text{ even}.
\end{cases}
\]
As established in the algebraic sections above, every orbit point satisfies $F(r_i)=E$; in particular,
\[
r_i\notin F,
\qquad\text{hence}\qquad
r_i\notin\{0,1\}
\]
for every round $i$.

We retain the reduced-descriptor notation introduced earlier for the evaluator
\[
T_{\tau,n}:F^N\to F^4.
\]
In this notation, the first round is represented by the fixed branch vectors
\[
\nu^{(0)}=(0,0,-1,1)^\top,
\qquad
\nu^{(1)}=(0,0,1,0)^\top,
\]
and every later round $i\ge 2$ is represented by the branch matrices
\[
M_i^{(0)},
\qquad
M_i^{(1)}.
\]
For a bit string $b=(b_1,\dots,b_n)\in\{0,1\}^n$, define
\[
w_{\tau,\rho}(b):=
\begin{cases}
\rho^\top\nu^{(b_1)},& n=1,\\[0.5ex]
\rho^\top M_n^{(b_n)}\cdots M_2^{(b_2)}\nu^{(b_1)},& n\ge 2,
\end{cases}
\]
for every scalarization vector $\rho\in F^4$.

For $u=(u_1,\dots,u_n)\in F^n$, define the affine interpolants
\[
\overline\nu(u_1):=(1-u_1)\nu^{(0)}+u_1\nu^{(1)},
\]
and, for $i\ge 2$,
\[
\overline M_i(u_i):=(1-u_i)M_i^{(0)}+u_iM_i^{(1)}.
\]
Then define the associated multilinear weight polynomial by
\[
\widetilde w_{\tau,\rho}(u):=
\begin{cases}
\rho^\top\overline\nu(u_1),& n=1,\\[0.5ex]
\rho^\top\overline M_n(u_n)\cdots \overline M_2(u_2)\overline\nu(u_1),& n\ge 2.
\end{cases}
\]
Finally, identify every vector $x\in F^N$ with a function
\[
x:\{0,1\}^n\to F,
\qquad
b\longmapsto x_b,
\]
and let
\[
\widetilde x\in F[U_1,\dots,U_n]
\]
denote its multilinear extension.

\begin{lemma}[Descriptor interpolation]
\label{lem:it-descriptor-interpolation}
For every $b\in\{0,1\}^n$,
\[
\widetilde w_{\tau,\rho}(b)=w_{\tau,\rho}(b).
\]
Moreover, $\widetilde w_{\tau,\rho}$ is multilinear in each variable, and for every point $u\in F^n$ its value $\widetilde w_{\tau,\rho}(u)$ is computable from $(d^{\mathrm{fold}}_{\tau,n},\rho)$ in $O(n)$ base-field operations.
\end{lemma}

\begin{proof}
At a Boolean point $b\in\{0,1\}^n$, every affine factor reduces to the branch that the corresponding bit selects:
\[
\overline\nu(b_1)=\nu^{(b_1)},
\qquad
\overline M_i(b_i)=M_i^{(b_i)}\quad(i\ge 2).
\]
Substituting these identities into the definition of $\widetilde w_{\tau,\rho}$ gives exactly the branch-product formula for $w_{\tau,\rho}(b)$.

Multilinearity is immediate because each factor depends affinely on its own variable and is independent of the others. For evaluation cost, fix $u\in F^n$ and $\rho\in F^4$. Computing $\widetilde w_{\tau,\rho}(u)$ then reduces to multiplying $n-1$ constant-size matrices into a constant-size vector and taking a final dot product with $\rho$. Each step costs $O(1)$ field operations, giving \(O(n)\) in total.
\end{proof}

\begin{lemma}[Scalarization identity]
\label{lem:it-scalarization}
For every $x\in F^N$ and every $\rho\in F^4$,
\[
\langle \rho,T_{\tau,n}(x)\rangle
=
\sum_{b\in\{0,1\}^n} x_b\,w_{\tau,\rho}(b)
=
\sum_{b\in\{0,1\}^n} \widetilde x(b)\,\widetilde w_{\tau,\rho}(b).
\]
\end{lemma}

\begin{proof}
Write
\[
x=\sum_{b\in\{0,1\}^n} x_b e_b,
\]
where $e_b\in F^N$ denotes the standard basis vector indexed by $b$. By linearity of $T_{\tau,n}$,
\[
T_{\tau,n}(x)=\sum_{b\in\{0,1\}^n} x_b T_{\tau,n}(e_b).
\]
Applying $\rho^\top$ yields
\[
\langle \rho,T_{\tau,n}(x)\rangle
=
\sum_{b\in\{0,1\}^n} x_b\,\rho^\top T_{\tau,n}(e_b)
=
\sum_{b\in\{0,1\}^n} x_b\,w_{\tau,\rho}(b).
\]
The second equality follows because $\widetilde x(b)=x_b$ at Boolean points and because Lemma~\ref{lem:it-descriptor-interpolation} gives $\widetilde w_{\tau,\rho}(b)=w_{\tau,\rho}(b)$.
\end{proof}

\subsection{The deterministic-orbit transcript is completely unsound}

Fix $x\in F^N$ and $\rho\in F^4$, and define
\[
H_{\tau,\rho}(u):=\widetilde x(u)\,\widetilde w_{\tau,\rho}(u)
\in F[U_1,\dots,U_n].
\]
Because both factors are multilinear, $H_{\tau,\rho}$ has individual degree at most $2$ in every variable. Since its coefficients lie in $F\subseteq E$, it can equally be evaluated at points in $E^n$.

Consider the following deterministic-orbit protocol for the claim
\[
\sum_{b\in\{0,1\}^n} H_{\tau,\rho}(b)=\alpha_0.
\]
The verifier uses the public orbit points $r_1,\dots,r_n\in E$ as its round challenges. At round $i$, the prover sends a polynomial $g_i(T)\in E[T]$ of degree at most $2$ satisfying
\[
g_i(0)+g_i(1)=\alpha_{i-1},
\]
and the verifier sets
\[
\alpha_i:=g_i(r_i).
\]
At the end, the verifier accepts if and only if
\[
\alpha_n=H_{\tau,\rho}(r_1,\dots,r_n).
\]

\begin{theorem}[Complete unsoundness of the deterministic-orbit transcript]
\label{thm:it-deterministic-unsound}
For every false initial claim $\alpha_0\in E$, there exists a cheating prover that makes the deterministic-orbit protocol above accept with probability $1$.
\end{theorem}

\begin{proof}
Fix any sequence of target intermediate claims
\[
\alpha_1,\dots,\alpha_{n-1}\in E,
\qquad
\alpha_n:=H_{\tau,\rho}(r_1,\dots,r_n).
\]
It is enough to show that, for every round $i$, any pair $(\alpha_{i-1},\alpha_i)\in E^2$ can be achieved by some degree-$\le 2$ polynomial $g_i(T)$ satisfying
\[
g_i(0)+g_i(1)=\alpha_{i-1},
\qquad
g_i(r_i)=\alpha_i.
\]
Consider a polynomial of the form
\[
g_i(T)=a_iT^2+b_iT.
\]
The two required conditions amount to the linear system
\[
a_i+b_i=\alpha_{i-1},
\qquad
a_i r_i^2+b_i r_i=\alpha_i.
\]
Its determinant is
\[
\det\begin{pmatrix}1&1\\ r_i^2&r_i\end{pmatrix}=r_i-r_i^2=r_i(1-r_i),
\]
which is nonzero because $r_i\notin\{0,1\}$, so the system has a unique solution $(a_i,b_i)\in E^2$ and every round constraint can be met exactly.

By construction, the final claim is set to equal the value the verifier computes in its final check. Therefore every round check and the final check pass, even though the initial claim $\alpha_0$ was false. The verifier accepts with probability $1$.
\end{proof}

\begin{remark}
This pathology requires nothing from the quartic structure except that the reused challenges satisfy $r_i\notin\{0,1\}$. The quartic reciprocal orbit simply provides a concrete deterministic family of such non-Boolean challenges.
\end{remark}

\subsection{The idealized oracle-model protocol}

The positive result concerns a \emph{different} protocol with three key changes. First, the prover sends its claimed output before seeing $\rho$. Second, the sumcheck rounds use fresh independent challenges drawn from $F$. Third, the verifier requires only the exact value $\widetilde x(z)$ at the final evaluation point $z\in F^n$.

For a fixed input $x\in F^N$, the protocol is:
\begin{enumerate}[label=\arabic*.]
\item The prover outputs a claimed evaluator output
\[
y\in F^4.
\]
This first message is sent before the verifier reveals $\rho$.
\item The verifier samples
\[
\rho\leftarrow F^4
\]
uniformly at random and sets
\[
\beta_0:=\langle \rho,y\rangle.
\]
\item The verifier defines
\[
H_{\tau,\rho}(u):=\widetilde x(u)\,\widetilde w_{\tau,\rho}(u).
\]
\item The prover and verifier run the standard $n$-round sumcheck protocol for the claim
\[
\sum_{b\in\{0,1\}^n} H_{\tau,\rho}(b)=\beta_0.
\]
At round $i$ the prover, who has seen $\rho$ and the earlier challenges $z_1,\dots,z_{i-1}$, sends a polynomial $h_i(T)\in F[T]$. The verifier checks that $\deg h_i\le 2$ and that
\[
h_1(0)+h_1(1)=\beta_0
\]
at the first round, and
\[
h_i(0)+h_i(1)=h_{i-1}(z_{i-1})
\]
for each later round. The verifier then samples
\[
z_i\leftarrow F
\]
uniformly and independently.
\item At the final point $z=(z_1,\dots,z_n)\in F^n$, the verifier queries the exact value $\widetilde x(z)$ and checks
\[
h_n(z_n)=H_{\tau,\rho}(z)=\widetilde x(z)\,\widetilde w_{\tau,\rho}(z).
\]
\end{enumerate}
Concretely, the prover transmits $h_i$ as its three coefficients in $F$, so the degree bound $\deg h_i\le 2$ is enforced syntactically.

\begin{lemma}[Standard degree-$2$ sumcheck soundness]
\label{lem:it-sumcheck}
Let
\[
H\in F[U_1,\dots,U_n]
\]
be any polynomial of individual degree at most $2$ in every variable. Consider the standard $n$-round sumcheck protocol for the claim
\[
\sum_{b\in\{0,1\}^n} H(b)=\beta_0,
\]
with verifier checks $\deg h_i\le 2$ at every round and with oracle access to $H(z_1,\dots,z_n)$ at the final evaluation point. Then every adaptive prover whose initial claim satisfies
\[
\beta_0\neq\sum_{b\in\{0,1\}^n} H(b)
\]
is accepted with probability at most $2n/q$.
\end{lemma}

\begin{proof}
We argue by induction on $n$.

For $n=1$, the verifier checks
\[
h_1(0)+h_1(1)=\beta_0
\]
and finally checks $h_1(z_1)=H(z_1)$. Since the initial claim is false,
\[
\beta_0\neq H(0)+H(1).
\]
Therefore any first-round polynomial $h_1$ satisfying the verifier's consistency equation must differ from $H$. Hence
\[
q_1(T):=h_1(T)-H(T)
\]
is a nonzero polynomial of degree at most $2$. Acceptance implies $q_1(z_1)=0$, which occurs with probability at most $2/q$.

Assume now that the statement holds for $(n-1)$ variables, and consider an $n$-variable instance. Define the true first-round partial-sum polynomial
\[
p_1(T):=\sum_{b_2,\dots,b_n\in\{0,1\}^{n-1}} H(T,b_2,\dots,b_n).
\]
Because $H$ has individual degree at most $2$, the univariate polynomial $p_1$ has degree at most $2$. Since the initial claim is false,
\[
p_1(0)+p_1(1)=\sum_{b\in\{0,1\}^n} H(b)\neq\beta_0.
\]
Thus any first-round prover message $h_1$ satisfying
\[
h_1(0)+h_1(1)=\beta_0
\]
must differ from $p_1$. Let
\[
q_1(T):=h_1(T)-p_1(T),
\]
so $q_1$ is a nonzero polynomial of degree at most $2$.

The verifier next samples $z_1\leftarrow F$. If $q_1(z_1)=0$, then the residual claim may become true; in this case we pessimistically bound the conditional acceptance probability by $1$. Since $q_1$ has at most $2$ roots,
\[
\Pr[q_1(z_1)=0]\le \frac{2}{q}.
\]
On the complementary event $q_1(z_1)\neq 0$, the residual claim for the specialized polynomial
\[
H(z_1,U_2,\dots,U_n)
\]
is still false, because the prover's carried-forward value is $h_1(z_1)$ while the true partial sum evaluates to $p_1(z_1)$. The continuation is therefore a standard $(n-1)$-round sumcheck instance, and its target polynomial still has individual degree at most $2$. By the induction hypothesis, conditioned on $q_1(z_1)\neq 0$, the remaining acceptance probability is at most $2(n-1)/q$.

Combining the two cases gives
\[
\Pr[\text{accept}]
\le
\Pr[q_1(z_1)=0]\cdot 1
+
\Pr[q_1(z_1)\neq 0]\cdot \frac{2(n-1)}{q}
\le
\frac{2}{q}+\frac{2(n-1)}{q}
=
\frac{2n}{q}.
\]
This completes the induction.
\end{proof}

\begin{proposition}[Why $y$ must be fixed before scalarization]
\label{prop:it-pre-rho-needed}
If the prover is allowed to choose $y$ after seeing $\rho\in F^4$, then the scalarization check alone provides no guarantees: for every $x\in F^N$ and every $\rho\in F^4$, there exists a vector
\[
y\neq T_{\tau,n}(x)
\]
such that
\[
\langle \rho,y\rangle=\langle \rho,T_{\tau,n}(x)\rangle.
\]
\end{proposition}

\begin{proof}
Choose any nonzero vector
\[
\delta\in\{v\in F^4 : \langle \rho,v\rangle=0\}.
\]
Such a nonzero vector always exists: if $\rho=0$, the orthogonal subspace is all of $F^4$, and if $\rho\neq 0$, the orthogonal complement $\rho^\perp$ has dimension $3$. Set
\[
y:=T_{\tau,n}(x)+\delta.
\]
Then $y\neq T_{\tau,n}(x)$ but
\[
\langle \rho,y\rangle
=
\langle \rho,T_{\tau,n}(x)\rangle+\langle \rho,\delta\rangle
=
\langle \rho,T_{\tau,n}(x)\rangle.
\]
Therefore no one-shot scalarization test can certify the vector equality $y=T_{\tau,n}(x)$ unless $y$ is fixed before $\rho$ is sampled.
\end{proof}

\begin{proposition}[Fresh scalarization alone does not repair the deterministic-orbit protocol]
\label{prop:it-rho-alone-not-enough}
Suppose the verifier samples $\rho\leftarrow F^4$ uniformly at random but still reuses the public orbit points $r_1,\dots,r_n$ as its round challenges. Then the resulting protocol remains completely unsound.
\end{proposition}

\begin{proof}
Fix any false output claim $y\neq T_{\tau,n}(x)$, and let
\[
\Delta:=y-T_{\tau,n}(x)\neq 0.
\]
For each realized value of $\rho$, there are two cases.

If $\langle \rho,\Delta\rangle=0$, then the scalarized initial claim
\[
\langle \rho,y\rangle=\langle \rho,T_{\tau,n}(x)\rangle
\]
is actually true. In this case the deterministic-orbit transcript passes all verifier checks.

If instead $\langle \rho,\Delta\rangle\neq 0$, then the scalarized initial claim is false. But Theorem~\ref{thm:it-deterministic-unsound} shows that, for the deterministic-orbit protocol, the prover can still force acceptance with probability $1$.

Thus for every realization of $\rho$ there exists an accepting prover strategy, so fresh scalarization without fresh verifier challenges does not restore soundness.
\end{proof}

\begin{theorem}[Information-theoretic oracle-model one-shot soundness]
\label{thm:it-oracle-soundness}
Consider the idealized protocol above. Let $P^\star$ be any possibly randomized prover that, on input $x\in F^N$, first outputs a vector $y\in F^4$ before $\rho$ is sampled, and then responds adaptively in the sumcheck phase, seeing each verifier challenge $z_j$ in turn. Then
\[
\Pr\big[\text{verifier accepts and }y\neq T_{\tau,n}(x)\big]
\le
\min\!\left\{1,\frac{2n+1}{q}\right\}.
\]
Here the probability is over $P^\star$'s internal coins and the uniform verifier randomness $(\rho,z_1,\dots,z_n)$.
\end{theorem}

\begin{proof}
Let $\omega$ denote the internal coins of $P^\star$, and fix an arbitrary realization of $\omega$. With $\omega$ fixed, the prover is deterministic. Let $y=y(\omega)$ be its first message, and define
\[
\Delta:=y-T_{\tau,n}(x)\in F^4.
\]
If $\Delta=0$, then the bad event
\[
\{\text{verifier accepts and } y\neq T_{\tau,n}(x)\}
\]
is impossible, so there is nothing to prove. Assume from now on that $\Delta\neq 0$.

Because $y$ is determined before $\rho$ is drawn, $\Delta$ does not depend on $\rho$. For nonzero $\Delta$, the map
\[
\rho\longmapsto \langle \rho,\Delta\rangle
\]
is a nonzero linear functional on the $4$-dimensional $F$-vector space $F^4$, and therefore its kernel has cardinality $q^3$. Consequently,
\[
\Pr_{\rho}\big[\langle \rho,\Delta\rangle=0\big]=\frac{1}{q}.
\]
Let
\[
\mathsf{Coll}:=\{\langle \rho,\Delta\rangle=0\}.
\]
On the complementary event $\neg\mathsf{Coll}$, we have
\[
\beta_0=\langle \rho,y\rangle\neq \langle \rho,T_{\tau,n}(x)\rangle.
\]
By Lemma~\ref{lem:it-scalarization},
\[
\langle \rho,T_{\tau,n}(x)\rangle
=
\sum_{b\in\{0,1\}^n} \widetilde x(b)\,\widetilde w_{\tau,\rho}(b)
=
\sum_{b\in\{0,1\}^n} H_{\tau,\rho}(b).
\]
Thus, for any fixed $\rho$ outside $\mathsf{Coll}$, the claimed sum $\beta_0$ differs from the true sum of $H_{\tau,\rho}$, so the prover faces a false sumcheck instance.

Fix any $\rho$ outside $\mathsf{Coll}$. By construction, both $\widetilde x$ and $\widetilde w_{\tau,\rho}$ are multilinear, so their product $H_{\tau,\rho}$ has individual degree at most $2$ in each variable. Therefore Lemma~\ref{lem:it-sumcheck} applies and gives
\[
\Pr[\text{accept}\mid \omega,\rho,\neg\mathsf{Coll}]\le \frac{2n}{q}.
\]
Bounding the acceptance probability by $1$ on the event $\mathsf{Coll}$, we obtain
\[
\Pr[\text{accept and }\Delta\neq 0\mid \omega]
\le
\Pr[\mathsf{Coll}\mid \omega]
+
\Pr[\text{accept}\mid \omega,\neg\mathsf{Coll}]
\le
\frac{1}{q}+\frac{2n}{q}
=
\frac{2n+1}{q}.
\]
Since the same bound holds for every fixed value of $\omega$, averaging over $\omega$ preserves the inequality. Finally, taking the minimum with $1$ gives the stated bound.
\end{proof}

\begin{remark}[Interpreting the error term]
The bound is nontrivial whenever $q\gg n$, which is the regime of practical interest. The two summands reflect distinct sources of error: the additive $1/q$ term is the unavoidable collision probability of the one-shot random scalarization map
\[
v\longmapsto \langle \rho,v\rangle,
\]
while the remaining $2n/q$ term comes from the $n$ rounds of degree-$2$ sumcheck.
\end{remark}

\begin{remark}[Exact scope of the theorem]
\label{rem:it-scope}
Theorem~\ref{thm:it-oracle-soundness} is a \emph{one-shot}, information-theoretic soundness statement in the oracle model: it assumes that the prover outputs $y$ before $\rho$ is sampled and bounds the error for a single evaluator-relation check $y=T_{\tau,n}(x)$.

Extensions to concrete polynomial commitments, adaptive multi-query opening, and Fiat--Shamir or correlated-challenge non-interactive compilation are separate layers that build on this core guarantee.
\end{remark}

\begin{thebibliography}{99}

\bibitem{BalbasCatalanoFioreLai2022CFC}
David Balb\'as, Dario Catalano, Dario Fiore, and Russell W.~F. Lai.
\newblock Chainable Functional Commitments for Unbounded-Depth Circuits.
\newblock \emph{IACR Cryptology ePrint Archive}, Paper 2022/1365, 2022.

\bibitem{BagadDombThaler2024SmallChar}
Suyash Bagad, Yuval Domb, and Justin Thaler.
\newblock The Sum-Check Protocol over Fields of Small Characteristic.
\newblock \emph{IACR Cryptology ePrint Archive}, Paper 2024/1046, 2024.

\bibitem{BlockEtAl2023FRI}
Alexander R. Block, Albert Garreta, Jonathan Katz, Justin Thaler, Pratyush Ranjan Tiwari, and Micha\l{} Zaj\k{a}c.
\newblock Fiat-Shamir Security of {FRI} and Related {SNARKs}.
\newblock \emph{IACR Cryptology ePrint Archive}, Paper 2023/1071, 2023.

\bibitem{BootleChiesaSotiraki2021}
Jonathan Bootle, Alessandro Chiesa, and Katerina Sotiraki.
\newblock Sumcheck Arguments and their Applications.
\newblock \emph{IACR Cryptology ePrint Archive}, Paper 2021/333, 2021.

\bibitem{CampanelliEtAl2022LMVC}
Matteo Campanelli, Anca Nitulescu, Carla R\`afols, Alexandros Zacharakis, and Arantxa Zapico.
\newblock Linear-Map Vector Commitments and their Practical Applications.
\newblock In \emph{ASIACRYPT 2022}. Full version: \emph{IACR Cryptology ePrint Archive}, Paper 2022/705.

\bibitem{BuenzChen2023ProtoStar}
Benedikt B\"unz and Binyi Chen.
\newblock ProtoStar: Generic Efficient Accumulation/Folding for Special Sound Protocols.
\newblock \emph{IACR Cryptology ePrint Archive}, Paper 2023/620, 2023.

\bibitem{CanettiEtAl2018FS}
Ran Canetti, Yilei Chen, Justin Holmgren, Alex Lombardi, Guy N. Rothblum, and Ron D. Rothblum.
\newblock Fiat-Shamir From Simpler Assumptions.
\newblock \emph{IACR Cryptology ePrint Archive}, Paper 2018/1004, 2018.

\bibitem{CatalanoFioreTucker2022AHFC}
Dario Catalano, Dario Fiore, and Ida Tucker.
\newblock Additive-Homomorphic Functional Commitments and Applications to Homomorphic Signatures.
\newblock In \emph{ASIACRYPT 2022}. Full version: \emph{IACR Cryptology ePrint Archive}, Paper 2022/1331.

\bibitem{FischLiuVesely2024Orbweaver}
Ben Fisch, Zeyu Liu, and Psi Vesely.
\newblock Orbweaver: Succinct Linear Functional Commitments from Lattices.
\newblock \emph{IACR Cryptology ePrint Archive}, Paper 2024/2026, 2024.

\bibitem{KothapalliSetty2023HyperNova}
Abhiram Kothapalli and Srinath Setty.
\newblock HyperNova: Recursive Arguments for Customizable Constraint Systems.
\newblock \emph{IACR Cryptology ePrint Archive}, Paper 2023/573, 2023.

\bibitem{KothapalliSettyTzialla2021Nova}
Abhiram Kothapalli, Srinath Setty, and Ioanna Tzialla.
\newblock Nova: Recursive Zero-Knowledge Arguments from Folding Schemes.
\newblock \emph{IACR Cryptology ePrint Archive}, Paper 2021/370, 2021.

\bibitem{LFKN92}
Carsten Lund, Lance Fortnow, Howard Karloff, and Noam Nisan.
\newblock Algebraic Methods for Interactive Proof Systems.
\newblock \emph{Journal of the ACM}, 39(4):859--868, 1992.

\bibitem{Rothblum2024MLENote}
Ron D. Rothblum.
\newblock A Note on Efficient Computation of the Multilinear Extension.
\newblock \emph{IACR Cryptology ePrint Archive}, Paper 2024/1103, 2024.

\end{thebibliography}

\end{document}
